% Results: Prolific Experiment
% Converted from prolific_results_prose_draft.md

\subsection{Sample and Exclusions}

We recruited 471 responses through Prolific. Table~\ref{tab:exclusion} reports the pre-registered exclusion pipeline. After removing survey previews (3), pre-registered pilot responses (17), and matching to Prolific's approved submissions, 446 responses remained. We then applied two pre-registered exclusion criteria: completion time under 120 seconds removed 89 responses (19.9\%), and the dual failure criterion --- failing both the instructed attention check and the comprehension check --- removed an additional 3 responses. The final analytical sample comprised 354 participants (Condition A: $n = 115$; Condition B: $n = 127$; Condition C: $n = 112$).

\begin{table}[htbp]
\centering
\small
\caption{Sample exclusion pipeline}
\label{tab:exclusion}
\begin{tabular}{llrl}
\toprule
Stage & $N$ Remaining & $N$ Excluded & Reason \\
\midrule
Qualtrics total        & 471 & ---  & Raw export \\
Remove previews        & 468 & 3   & Survey preview responses \\
Remove pilots          & 451 & 17  & Pre-registered pilot exclusion \\
Prolific matching      & 446 & 5   & Unmatched to Prolific approved submissions \\
Speed exclusion        & 357 & 89  & Completion time $<$ 120 seconds \\
Dual-check exclusion   & 354 & 3   & Failed both attention and comprehension checks \\
\bottomrule
\end{tabular}
\end{table}


The speed exclusion rate warrants comment. Approximately one in five participants completed the survey in under two minutes, a duration insufficient to read the vignette, process the performance notification, and respond thoughtfully to the dependent measures. Among the 357 participants who met the speed threshold, only 3 failed the dual attention-comprehension criterion, suggesting that the speed filter effectively separated engaged from disengaged respondents.

A condition balance check revealed a significant age imbalance: participants in Condition C were older ($M = 41.2$) than those in Condition A ($M = 37.0$) and Condition B ($M = 37.5$), $F(2, 351) = 4.15$, $p = .017$. All other covariates --- financial literacy, gender, education, income, and robo-advisor experience --- were balanced across conditions. Age is included as a covariate in all regression models.

\subsection{Descriptive Statistics}

\begin{table}[htbp]
\centering
\small
\caption{Means and standard deviations of dependent variables by experimental condition}
\label{tab:descriptives}
\begin{tabular}{lcccc}
\toprule
 & Condition A & Condition B & Condition C & Overall \\
 & ($n = 115$) & ($n = 127$) & ($n = 112$) & ($N = 354$) \\
\midrule
Evaluative trust ($\alpha = .43$) & 5.43 (1.01) & 5.28 (1.12) & 5.19 (1.01) & 5.30 (1.05) \\
Delegative trust ($\alpha = .81$) & 3.84 (1.39) & 4.04 (1.48) & 3.86 (1.39) & 3.92 (1.42) \\
Betrayal                          & 4.23 (1.61) & 4.02 (1.71) & 4.24 (1.61) & 4.16 (1.64) \\
Disappointment                    & 5.84 (1.16) & 5.71 (1.36) & 5.71 (1.18) & 5.75 (1.24) \\
Betrayal premium                  & $-1.62$ (1.57) & $-1.69$ (1.65) & $-1.46$ (1.63) & $-1.59$ (1.62) \\
General exit                      & 4.68 (1.60) & 4.57 (1.72) & 4.76 (1.58) & 4.67 (1.64) \\
Algorithm aversion                & 5.03 (1.53) & 4.65 (1.58) & 4.95 (1.43) & 4.86 (1.52) \\
\bottomrule
\end{tabular}
\begin{flushleft}
\footnotesize
\textit{Note.} All items measured on 7-point Likert scales (1 = Strongly Disagree, 7 = Strongly Agree). Betrayal premium = Betrayal minus Disappointment. Condition A = Fiduciary/Relationship framing; Condition B = Performance/Analytics framing; Condition C = Neutral control. Cronbach's $\alpha$ reported for composite measures.
\end{flushleft}
\end{table}


Table~\ref{tab:descriptives} presents means and standard deviations for all dependent variables by condition. Two patterns merit attention. First, evaluative trust ($M = 5.30$, $SD = 1.05$) was substantially higher than delegative trust ($M = 3.92$, $SD = 1.42$) across all conditions. Participants were generally willing to monitor and evaluate the algorithm's performance but ambivalent about granting it discretionary authority --- a pattern consistent with the framework's prediction that the default psychological posture toward algorithms is evaluative. Second, disappointment ($M = 5.75$) exceeded betrayal ($M = 4.16$) in every condition, producing a consistently negative betrayal premium ($M = -1.59$). After an adverse event, participants' dominant reactive attitude was disappointment at the outcome rather than betrayal by the agent, again consistent with evaluative rather than delegative trust governing the human--algorithm relationship.

The delegative trust composite showed good internal consistency (Cronbach's $\alpha = .807$). The evaluative trust composite, comprising only two items, showed poor internal consistency ($\alpha = .429$). We discuss implications of this measurement limitation below and in Section~\ref{sec:discussion}.

Table~\ref{tab:correlations} presents bivariate correlations among the trust composites and dependent variables. The evaluative and delegative trust composites are positively correlated ($r = .246$, $p < .001$), confirming that they are related constructs --- as any two trust measures should be --- while sharing only approximately 6\% of variance. The two composites thus capture substantially distinct dimensions of the participant's orientation toward the algorithm. Their downstream profiles reinforce this conclusion. Delegative trust is negatively correlated with algorithm aversion ($r = -.313$, $p < .001$) and general exit ($r = -.155$, $p < .01$), but uncorrelated with betrayal ($r = .061$, n.s.). Evaluative trust shows a largely flat profile across dependent variables, consistent with both its lower reliability and the descriptive observation that evaluative engagement is near-uniformly high across the sample. Among the post-shock dependent variables, general exit and algorithm aversion are strongly correlated ($r = .690$, $p < .001$), reflecting their shared orientation toward abandoning the algorithm, though they differ in specificity: general exit captures platform-switching intent while algorithm aversion captures the preference for human advice specifically.

\begin{table}[htbp]
\centering
\small
\caption{Bivariate correlations among trust composites and dependent variables}
\label{tab:correlations}
\begin{tabular}{lccccccc}
\toprule
 & 1 & 2 & 3 & 4 & 5 & 6 & 7 \\
\midrule
1. Evaluative trust  & ---    &        &        &        &        &        &     \\
2. Delegative trust  & .246*** & ---    &        &        &        &        &     \\
3. Betrayal          & .054   & .061   & ---    &        &        &        &     \\
4. Disappointment    & .113*  & $-.102$ & .399*** & ---    &        &        &     \\
5. Betrayal premium  & $-.032$ & .140** & .712*** & $-.359$*** & ---    &        &     \\
6. General exit      & $-.015$ & $-.155$** & .578*** & .406*** & .278*** & ---    &     \\
7. Algorithm aversion & $-.109$* & $-.313$*** & .515*** & .420*** & .202*** & .690*** & --- \\
\bottomrule
\end{tabular}
\begin{flushleft}
\footnotesize
\textit{Note.} $N = 354$. Evaluative trust = mean of Items 1--2; Delegative trust = mean of Items 3--5; Betrayal premium = Betrayal minus Disappointment. \\
* $p < .05$, ** $p < .01$, *** $p < .001$.
\end{flushleft}
\end{table}


\subsection{Experimental Manipulation}

We predicted that framing the robo-advisor in fiduciary terms (Condition A) would produce higher delegative trust than framing it in performance terms (Condition B), and that the performance framing would produce higher evaluative trust (H1). Neither prediction was supported. One-way ANOVA revealed no significant differences across conditions in delegative trust, $F(2, 351) = 0.75$, $p = .473$, $\eta^2 = .004$, or evaluative trust, $F(2, 351) = 1.50$, $p = .225$, $\eta^2 = .009$ (Table~\ref{tab:anova}). The observed means were directionally opposite to predictions on both measures, though far from significance.

\begin{table}[htbp]
\centering
\small
\caption{One-way ANOVA results for experimental manipulation effects (H1, H2)}
\label{tab:anova}
\begin{tabular}{llccccl}
\toprule
Hypothesis & DV & $F(2, 351)$ & $p$ & $\eta^2$ & Predicted & Observed \\
\midrule
H1a & Delegative trust  & 0.75 & .473 & .004 & A $>$ B & B $>$ A (n.s.) \\
H1b & Evaluative trust  & 1.50 & .225 & .009 & B $>$ A & A $>$ B (n.s.) \\
H2  & Betrayal premium  & 0.57 & .564 & .003 & A $>$ B & A $>$ B (n.s.) \\
\bottomrule
\end{tabular}
\begin{flushleft}
\footnotesize
\textit{Note.} Condition A = Fiduciary/Relationship framing; Condition B = Performance/Analytics framing; Condition C = Neutral control. n.s.\ = not significant at $\alpha = .05$. $N = 354$.
\end{flushleft}
\end{table}


The betrayal premium (H2) was likewise undifferentiated across conditions, $F(2, 351) = 0.57$, $p = .564$, $\eta^2 = .003$. All three conditions produced negative betrayal premiums, indicating that disappointment dominated betrayal regardless of framing.

We interpret the manipulation failure in Section~\ref{sec:discussion}. Crucially, the null effect of condition assignment does not bear on the validity of the trust constructs as individual-difference measures. The primary predictions of the evaluative--delegative framework concern the downstream consequences of trust modes, not their experimental inducibility, and we test these predictions in the analyses that follow.

\subsection{Delegative Trust and Post-Shock Resilience}

Our framework predicts that delegative trust --- granting the advisor discretionary authority with acceptance of opacity --- should buffer against exit-seeking behavior after an adverse event, because the delegative truster's relationship with the advisor is not contingent on any particular performance outcome. H4 tests this prediction directly.

Table~\ref{tab:reg_h4} reports the full regression model predicting general exit intent. The model is significant, $F(9, 344) = 2.18$, $p = .023$, $R^2 = .054$. As predicted, delegative trust negatively predicts exit intent ($B = -0.206$, $SE = 0.064$, $\beta = -0.179$, $p = .001$, $95\%$ CI $[-0.33, -0.08]$). Each unit increase in delegative trust on the seven-point scale is associated with a 0.21-point decrease in the intention to research alternative platforms. Evaluative trust, by contrast, shows no predictive relationship with exit intent ($B = 0.037$, $p = .667$). The two trust dimensions thus make divergent predictions about post-shock behavior: one predicts resilience, the other is inert.

\begin{table}[htbp]
\centering
\small
\caption{OLS regression predicting general exit intent (H4)}
\label{tab:reg_h4}
\begin{tabular}{lcccccc}
\toprule
Predictor & $B$ & SE & $\beta$ & $t$ & $p$ & 95\% CI \\
\midrule
\textbf{Delegative trust}    & \textbf{$-0.206$**} & \textbf{0.064} & \textbf{$-0.179$} & \textbf{$-3.23$} & \textbf{.001} & \textbf{[$-0.33$, $-0.08$]} \\
Evaluative trust              & 0.037                & 0.087          & 0.024              & 0.43              & .667 & [$-0.13$, $0.21$] \\
Condition A (Fiduciary)       & $-0.056$             & 0.218          & $-0.016$           & $-0.26$           & .797 & [$-0.48$, $0.37$] \\
Condition B (Performance)     & $-0.132$             & 0.211          & $-0.039$           & $-0.62$           & .532 & [$-0.55$, $0.28$] \\
Age                           & 0.005                & 0.007          & 0.040              & 0.74              & .460 & [$-0.01$, $0.02$] \\
Financial literacy            & $-0.251$             & 0.134          & $-0.106$           & $-1.87$           & .062 & [$-0.51$, $0.01$] \\
Robo-advisor experience       & 0.242                & 0.266          & 0.049              & 0.91              & .365 & [$-0.28$, $0.76$] \\
Comprehension pass            & 0.616**              & 0.232          & 0.146              & 2.65              & .008 & [$0.16$, $1.07$] \\
Attention pass                & 0.577                & 0.478          & 0.064              & 1.21              & .228 & [$-0.36$, $1.52$] \\
\midrule
\multicolumn{7}{l}{$R^2 = .054$, Adjusted $R^2 = .029$, $F(9, 344) = 2.18$, $p = .023$} \\
\bottomrule
\end{tabular}
\begin{flushleft}
\footnotesize
\textit{Note.} Reference category for condition is C (Neutral control). Robo-advisor experience coded as 1 = has used a robo-advisor. Comprehension pass = correctly identified WealthPath underperformance. Attention pass = selected instructed response. $N = 354$. \\
* $p < .05$, ** $p < .01$, *** $p < .001$.
\end{flushleft}
\end{table}


Two additional findings are worth noting. Comprehension check pass positively predicts exit intent ($B = 0.616$, $p = .008$): participants who correctly identified the performance shortfall were more likely to seek alternatives. This pattern is consistent with evaluative processing of the adverse event --- the negative signal was encoded and acted upon --- though the effect is captured by the check variable rather than the evaluative trust composite, a point to which we return below. Financial literacy approaches significance as a negative predictor ($B = -0.251$, $p = .062$), suggesting that higher-literacy participants may be more tolerant of short-term underperformance.

The result is robust to alternative exclusion criteria. Under strict attention exclusion ($N = 342$), the delegative trust coefficient is $B = -0.207$, $p = .001$. Under strict comprehension exclusion ($N = 289$), it is $B = -0.161$, $p = .021$. Direction and significance are unchanged.

\subsection{Delegative Trust and Algorithm Aversion}

H3 tested whether the trust dimensions predict algorithm aversion --- the preference for a human advisor or self-management over the algorithm after the adverse event. We predicted that evaluative trust would drive algorithm aversion and that delegative trust would not. The results reversed this prediction in a theoretically informative way.

Table~\ref{tab:reg_h3} reports the full model, which is highly significant, $F(9, 344) = 6.82$, $p < .001$, $R^2 = .151$, and represents the strongest overall fit of any regression in the experimental analyses. Delegative trust is the dominant predictor of algorithm aversion ($B = -0.326$, $SE = 0.056$, $\beta = -0.304$, $p < .001$, $95\%$ CI $[-0.44, -0.22]$). Participants who scored higher on willingness to grant the algorithm discretionary authority were substantially less likely to prefer a human advisor after the adverse event. Evaluative trust, by contrast, was not significant ($B = -0.039$, $p = .614$).

\begin{table}[htbp]
\centering
\small
\caption{OLS regression predicting algorithm aversion (H3)}
\label{tab:reg_h3}
\begin{tabular}{lcccccc}
\toprule
Predictor & $B$ & SE & $\beta$ & $t$ & $p$ & 95\% CI \\
\midrule
\textbf{Delegative trust}    & \textbf{$-0.326$***} & \textbf{0.056} & \textbf{$-0.304$} & \textbf{$-5.79$} & \textbf{$< .001$} & \textbf{[$-0.44$, $-0.22$]} \\
Evaluative trust              & $-0.039$             & 0.077          & $-0.027$           & $-0.50$           & .614 & [$-0.19$, $0.11$] \\
Condition A (Fiduciary)       & 0.073                & 0.192          & 0.023              & 0.38              & .703 & [$-0.30$, $0.45$] \\
Condition B (Performance)     & $-0.250$             & 0.186          & $-0.079$           & $-1.35$           & .179 & [$-0.62$, $0.12$] \\
Age                           & $-0.006$             & 0.006          & $-0.051$           & $-1.01$           & .315 & [$-0.02$, $0.01$] \\
Financial literacy            & $-0.191$             & 0.118          & $-0.086$           & $-1.62$           & .107 & [$-0.42$, $0.04$] \\
Robo-advisor experience       & $-0.215$             & 0.234          & $-0.047$           & $-0.92$           & .358 & [$-0.68$, $0.25$] \\
Comprehension pass            & 0.584**              & 0.204          & 0.149              & 2.86              & .004 & [$0.18$, $0.99$] \\
Attention pass                & 1.179**              & 0.421          & 0.141              & 2.80              & .005 & [$0.35$, $2.01$] \\
\midrule
\multicolumn{7}{l}{$R^2 = .151$, Adjusted $R^2 = .129$, $F(9, 344) = 6.82$, $p < .001$} \\
\bottomrule
\end{tabular}
\begin{flushleft}
\footnotesize
\textit{Note.} Reference category for condition is C (Neutral control). Robo-advisor experience coded as 1 = has used a robo-advisor. Comprehension pass = correctly identified WealthPath underperformance. Attention pass = selected instructed response. $N = 354$. \\
* $p < .05$, ** $p < .01$, *** $p < .001$.
\end{flushleft}
\end{table}


This finding has a straightforward interpretation within the evaluative--delegative framework. The factor that protects against algorithm aversion is not a favorable performance assessment --- it is the willingness to cede decisional authority. Participants who had granted the algorithm discretion, who had accepted opacity and declined to second-guess every decision, maintained that posture through an adverse event. Those who had not granted such discretion showed no corresponding resilience. This is dimensional separability operating in the predicted direction: the two trust modes engage different psychological mechanisms, and only the delegative mechanism predicts who stays with the algorithm after failure.

As in H4, the check variables are informative. Both comprehension pass ($B = 0.584$, $p = .004$) and attention pass ($B = 1.179$, $p = .005$) positively predict algorithm aversion. Participants who were paying attention and who correctly encoded the performance shortfall were more, not less, likely to want a human advisor. This pattern is consistent with evaluative processing --- the participants who received and encoded the negative signal were the ones who responded by seeking alternatives --- but the evaluative trust composite does not capture this effect. The poor internal consistency of the evaluative trust composite ($\alpha = .429$) likely explains the discrepancy: the comprehension and attention check variables are acting as proxies for evaluative engagement that the two-item scale measures inadequately. This is a limitation of the present study; we discuss it further in Section~\ref{sec:discussion}, but note here that the limitation biases against finding evaluative trust effects and does not threaten the delegative trust findings.

This result also held under both robustness checks. Strict attention exclusion ($N = 342$): delegative trust $B = -0.322$, evaluative trust $B = -0.058$ (n.s.). Strict comprehension exclusion ($N = 289$): delegative trust $B = -0.316$, evaluative trust $B = -0.047$ (n.s.).

\subsection{Financial Literacy Interaction}

H5 predicted that financial literacy would moderate the effect of condition on delegative trust, with high-literacy participants showing a larger difference between Conditions A and B. The interaction was not significant (Condition A $\times$ high literacy: $B = 0.072$, $p = .869$; Condition B $\times$ high literacy: $B = -0.109$, $p = .797$). This null follows from the manipulation failure: where condition assignment does not produce main effects on delegative trust, an interaction with literacy is unlikely to emerge.

One finding from the H5 model merits attention independent of the interaction. Prior robo-advisor experience was the strongest predictor of delegative trust in the model ($B = 0.729$, $p = .001$). Participants who had previously used a robo-advisor scored nearly three-quarters of a point higher on the delegative trust composite, regardless of experimental condition or financial literacy. This is a substantial effect that dwarfs any condition-level variation and is consistent with the philosophical account of trust developed in Section~\ref{sec:theory}: trust postures are learned through repeated encounters rather than installed by informational framing. We return to this finding in Section~\ref{sec:discussion}.

\subsection{Exploratory Analyses}

Three pre-registered exploratory analyses are reported here. These were designated as exploratory in the pre-registration and should not be treated as confirmatory.

We predicted that Condition A would produce higher variance in delegative trust than the other conditions (E1), reflecting polarization between participants who can and cannot extend delegative trust to algorithms. This was not supported: Levene's test was not significant, $F(2, 351) = 0.32$, $p = .728$, and variances were similar across conditions (A: 1.94; B: 2.18; C: 1.92).

We predicted that comprehension check failures would show higher delegative trust and lower evaluative trust than passes (E2), on the theory that delegative trusters encode performance signals less carefully. The predicted directions emerged on both measures --- failures showed slightly higher delegative trust ($M = 3.97$ vs.\ $3.90$) and lower evaluative trust ($M = 5.11$ vs.\ $5.34$) --- but neither difference reached significance (delegative: $t = -0.33$, $p = .739$; evaluative: $t = 1.60$, $p = .110$). Comprehension failure rates did not differ meaningfully by condition (A: 19.1\%; B: 18.9\%; C: 17.0\%; Table~\ref{tab:comprehension}).

\begin{table}[htbp]
\centering
\small
\caption{Comprehension check performance by experimental condition}
\label{tab:comprehension}
\begin{tabular}{lccc}
\toprule
Condition & Pass $N$ (\%) & Fail $N$ (\%) & Total \\
\midrule
A (Fiduciary)    & 93 (80.9\%)  & 22 (19.1\%) & 115 \\
B (Performance)  & 103 (81.1\%) & 24 (18.9\%) & 127 \\
C (Neutral)      & 93 (83.0\%)  & 19 (17.0\%) & 112 \\
\midrule
Total            & 289 (81.6\%) & 65 (18.4\%) & 354 \\
\bottomrule
\end{tabular}
\begin{flushleft}
\footnotesize
\textit{Note.} Comprehension check asked participants to identify WealthPath's performance relative to the S\&P 500. Correct answer: WealthPath performed worse than the market.
\end{flushleft}
\end{table}


We tested whether prior robo-advisor experience moderated the relationship between condition and each dependent variable (E3). Two of fourteen interactions reached significance: Condition A $\times$ robo-experience on disappointment ($B = 1.174$, $p = .022$), and Condition B $\times$ robo-experience on the betrayal premium ($B = -1.313$, $p = .035$). Among participants with robo-advisor experience, fiduciary framing produced greater disappointment after the adverse event, while performance framing produced a more negative betrayal premium --- that is, a stronger dominance of disappointment over betrayal. Both patterns are consistent with the framework: experienced users who are told the algorithm is their fiduciary may have activated higher expectations, while experienced users primed with performance framing process the shortfall as a pure evaluative failure. However, with fourteen interaction tests conducted at $\alpha = .05$, approximately 0.7 significant results are expected by chance. These two interactions modestly exceed chance expectation but would not survive correction for multiple comparisons. We treat them as suggestive and hypothesis-generating.

\subsection{Summary}

The experimental manipulation did not differentiate trust modes across conditions, but the evaluative--delegative distinction proves consequential when measured as a continuous individual difference. The two trust composites share only 6\% of variance ($r = .246$) and show divergent downstream profiles (Table~\ref{tab:correlations}). Delegative trust is the primary predictor of resilience to both exit-seeking (H4: $\beta = -0.179$, $p = .001$) and algorithm aversion (H3: $\beta = -0.304$, $p < .001$) following an adverse event. Evaluative trust predicts neither outcome, though this null must be interpreted cautiously given the low reliability of the two-item composite. The trust dimensions thus make divergent and specific predictions about post-shock behavior, supporting the core theoretical claim of dimensional separability. These findings complement the CFPB analysis reported in Section~\ref{sec:results_cfpb}, which provides naturalistic evidence for the same constructs; we discuss the triangulation in Section~\ref{sec:discussion}.
